% H3 Hiemenz solution — Persian article, XeLaTeX + xepersian
\documentclass[12pt,a4paper]{article}
\usepackage{amsmath,amssymb}
\usepackage{booktabs}
\usepackage{graphicx}
\usepackage{geometry}
\geometry{margin=25mm}
\usepackage[hidelinks]{hyperref}
\usepackage{xepersian}
\settextfont{Vazirmatn}
\setdigitfont{Vazirmatn}
\usepackage{setspace}
\setstretch{1.25}
\begin{document}

\begin{center}
{\Large\bfseries حل جریان نقطهٔ سکون هایمنز\par}
\vspace{1ex}
{\large تبدیل تشابهی، استخراج معادله، و حل عددی\par}
\vspace{1ex}
{\small امیر مرادی\\
درس لایهٔ مرزی · ۶ آبان ۱۴۰۵ / \lr{2026-10-28}\par}
\end{center}

\begin{center}{\large\bfseries چکیده}\end{center}
جریان نقطهٔ سکون با فرض تشابهی به یک معادلهٔ دیفرانسیل معمولیِ مرتبهٔ سه
تبدیل می‌شود: \lr{$F''' + F F'' - (F')^2 + 1 = 0$} با
\lr{$F(0)=F'(0)=0$} و \lr{$F'(\infty)=1$}. حل عددیِ این مقاله،
\lr{$F''(0)=1.2325876568$} و ضخامت جابه‌جاییِ بی‌بعد
\lr{$\delta=0.64790047$} می‌دهد.

\noindent\textbf{واژگان کلیدی:} هایمنز؛ نقطهٔ سکون؛ لایهٔ مرزی؛ حل تشابهی؛
فالکنر--اسکن؛ تنش برشی دیوار.

\section{جواب اول}
با \lr{$\eta=x_2^*$}، سرعت‌های بی‌بعد چنین‌اند:
\lr{$u_1^*=x_1^*F'(\eta)$} و \lr{$u_2^*=-F(\eta)$}، و \lr{$F$} جوابِ
\lr{$F''' + F F'' - (F')^2 + 1 = 0$} با \lr{$F(0)=F'(0)=0$} و
\lr{$F'(\infty)=1$} است. مهم‌ترین عددِ حل، \lr{$F''(0)=1.2325876568$}
است که تنش برشی دیوار با آن متناسب است.

\section{صورت مسئله و متغیرهای بی‌بعد}
جریان پایا، دوبعدی و تراکم‌ناپذیر روی صفحهٔ تختِ \lr{$x_2=0$}، با جریان
بیرونیِ \lr{$U_e=ax_1$}، \lr{$a>0$}، و لزجت سینماتیکیِ \lr{$\gamma$}
در نظر گرفته می‌شود. طول مشخصه
\lr{$\sqrt{\gamma/a}$} و سرعت مشخصه \lr{$\sqrt{\gamma a}$}‌اند، پس:
\begin{equation}
x_1^*, x_2^* = \frac{x_1, x_2}{\sqrt{\gamma/a}},
\qquad
u_1^*, u_2^* = \frac{u_1, u_2}{\sqrt{\gamma a}}.
\end{equation}
در این مقیاس، معادلهٔ مومنتمِ لایهٔ مرزیِ بی‌بعد چنین می‌شود:
\begin{equation}
u_1^*\frac{\partial u_1^*}{\partial x_1^*}
+ u_2^*\frac{\partial u_1^*}{\partial x_2^*}
= x_1^* + \frac{\partial^2 u_1^*}{\partial {x_2^*}^2},
\label{eq:bl}
\end{equation}
که جملهٔ \lr{$x_1^*$} سمت راست، همان گرادیان فشارِ جریان بیرونیِ نقطهٔ
سکون است.

\section{پیوستگی و استخراج معادله}
با فرض \lr{$u_1^*=x_1^*F'(\eta)$} و \lr{$u_2^*=-F(\eta)$}،
\lr{$\eta=x_2^*$}:
\begin{equation}
\frac{\partial u_1^*}{\partial x_1^*} = F'(\eta),
\qquad
\frac{\partial u_2^*}{\partial x_2^*} = -F'(\eta),
\end{equation}
پس پیوستگی، \lr{$F'-F'=0$}، برقرار است. مشتق‌های لازم برای
\eqref{eq:bl}:
\begin{equation}
\frac{\partial u_1^*}{\partial x_2^*} = x_1^*F''(\eta),
\qquad
\frac{\partial^2 u_1^*}{\partial {x_2^*}^2} = x_1^*F'''(\eta).
\end{equation}
جای‌گذاری در \eqref{eq:bl}:
\begin{equation}
x_1^*(F')^2 - x_1^*FF'' = x_1^* + x_1^*F''',
\end{equation}
و با تقسیم بر \lr{$x_1^*$} و مرتب‌سازی:
\begin{equation}
\boxed{\,F''' + FF'' - (F')^2 + 1 = 0\,}.
\end{equation}
جملهٔ \lr{$+1$}، همان جملهٔ گرادیان فشار است؛ بدون آن، معادلهٔ بلازیوس
به دست می‌آید. به همین دلیل، این معادله، حالتِ \lr{$\beta=1$} معادلهٔ
فالکنر--اسکن است.

\section{شرایط مرزی}
\begin{equation}
F(0)=0, \qquad F'(0)=0, \qquad F'(\infty)=1.
\end{equation}
\lr{$F(0)=0$} یعنی \lr{$u_2^*=0$} روی دیوار: نفوذناپذیری.
\lr{$F'(0)=0$} یعنی \lr{$u_1^*=0$} روی دیوار: چسبندگی.
\lr{$F'(\infty)=1$} یعنی در دوردست، \lr{$u_1^*\to x_1^*$}، یعنی سرعت
به سرعت جریان بیرونیِ \lr{$U_e=ax_1$} می‌رسد. معادله مرتبهٔ سه است و
دقیقاً سه شرط دارد؛ پس مسئله خوش‌تعریف است و مجهولِ شوتینگ، فقط
\lr{$F''(0)$} است.

\section{روش عددی و اعتبارسنجی}
به دستور استاد، حل با \textbf{رانگ--کوتای مرتبهٔ چهار} و روش شوتینگ
انجام شد. کد، ساده و دانشجویی است و پلات ندارد: فقط با کتابخانهٔ
استاندارد پایتون، \lr{$F''(0)$} را حدس می‌زند، با رانگ--کوتا تا
\lr{$\eta=6$} جلو می‌رود، و با تنصیف در بازهٔ \lr{$[1.0,1.5]$} آن را
چنان تنظیم می‌کند که \lr{$F'(6)=1$} شود. تا \lr{$\eta=6$}، مقدار
\lr{$F'$} تا \lr{$10^{-8}$} به یک رسیده است، پس \lr{$\eta=6$} جانشین
خوبی برای بی‌نهایت است.

نکتهٔ عملی شوتینگ: برای حدس‌های دور از جواب، حلِ مقدار اولیه به‌صورت
نمایی منفجر می‌شود، پس بازهٔ تنصیف را نباید خیلی باز گرفت؛ بازهٔ
\lr{$[1.0,1.5]$} برای این مسئله کافی و پایدار است.

اعتبارسنجی (پالایش گام):
\begin{itemize}
\item \lr{$h=0.02$}: \lr{$F''(0)=1.2325876554$}.
\item \lr{$h=0.01$}: \lr{$F''(0)=1.2325876568$}.
\item \lr{$h=0.005$}: \lr{$F''(0)=1.2325876569$}.
\end{itemize}
پس با ریز شدن گام، جواب تا حدود \lr{$10^{-9}$} پایدار است.

\section{نتایج}

\begin{table}[htbp]
\centering
\caption{نتایج واقعی اجرا (فایل \lr{\texttt{results/tables/metrics.json}})}
\begin{tabular}{cccc}
\toprule
\lr{$\eta$} & \lr{$F$} & \lr{$F'$} & \lr{$F''$} \\
\midrule
0 & 0.00000000 & 0.00000000 & 1.23258766 \\
0.5 & 0.13358522 & 0.49464929 & 0.75830748 \\
1 & 0.45922702 & 0.77786527 & 0.39801295 \\
1.5 & 0.88732899 & 0.91616823 & 0.17695812 \\
2 & 1.36197416 & 0.97321674 & 0.06582538 \\
2.5 & 1.85442873 & 0.99285118 & 0.02022731 \\
3 & 2.35255668 & 0.99842416 & 0.00507796 \\
3.5 & 2.85217346 & 0.99971588 & 0.00103216 \\
4 & 3.35210930 & 0.99995843 & 0.00016867 \\
4.5 & 3.85210058 & 0.99999510 & 0.00002204 \\
5 & 4.35209962 & 0.99999954 & 0.00000229 \\
5.5 & 4.85209953 & 0.99999997 & 0.00000019 \\
6 & 5.35209953 & 1.00000000 & 0.00000001 \\
\bottomrule
\end{tabular}
\end{table}

در دوردست، \lr{$F\approx\eta-\delta$} با \lr{$\delta=0.64790047$}؛ پس
ضخامت جابه‌جایی، \lr{$\delta^*=0.6479\sqrt{\gamma/a}$} است. سرعت مماسی
در \lr{$\eta=1$} حدود ۷۸٪ سرعت بیرونی، در \lr{$\eta=2.5$}
حدود ۹۹٪ (\lr{$F'=0.99285118$})، و در \lr{$\eta=4$} عملاً برابر
آن است؛ پس ضخامت لایهٔ مرزی، چند برابرِ \lr{$\sqrt{\gamma/a}$} است.
تنش برشی دیوار:
\begin{equation}
\tau_w = \mu\left.\frac{\partial u_1}{\partial x_2}\right|_{x_2=0}
= \rho\sqrt{\gamma a^3}\,x_1\,F''(0),
\end{equation}
یعنی با \lr{$x_1$} خطی زیاد می‌شود و ضریبش همان \lr{$1.2325876568$}
است.

\section{منابع}
\begin{enumerate}
\item \lr{Hiemenz, K. (1911). Die Grenzschicht an einem in den gleichf\"ormigen Fl\"ussigkeitsstrom eingetauchten geraden Kreiszylinder. Dinglers Polytechnisches Journal, 326, 321--324.}
\item \lr{Schlichting, H., \& Gersten, K. Boundary-Layer Theory. Springer.}
\item \lr{White, F. M. Viscous Fluid Flow. McGraw-Hill.}
\item کد و نتایج این مقاله:
\lr{\texttt{notebooks/hiemenz\_rk4.py}} (رانگ--کوتا، بدون پلات)،
\lr{\texttt{results/tables/metrics.json}}.
\end{enumerate}

\end{document}
